\section{Related Work}
\label{sec:related}

We organize prior work by the type of signal used for rank selection, showing that all existing approaches require training or calibration data, and identify the structural approach as a novel alternative.

\subsection{Gradient-Based Adaptive Rank Methods}

The dominant approach to per-layer rank selection uses gradient signals during training to estimate layer importance.

\textbf{AdaLoRA} \cite{Zhang2023AdaLoRA} parameterizes weight updates in singular value decomposition form and prunes singular values based on importance scores derived from gradient magnitudes. Its rank allocation is dynamic, changing during training. While effective, AdaLoRA requires training through the adaptive phase and carries substantial optimizer overhead. Critically, the singular values being pruned belong to $\Delta W$ (the learned update), not $W_0$ (the pretrained weights) --- a fundamental architectural difference from our approach.

\textbf{DyLoRA} \cite{Valipour2022DyLoRA} trains the LoRA adapter simultaneously across a range of ranks, enabling rank selection post-training by truncating the adapter. This eliminates grid search but still requires full training before the optimal rank is known.

\textbf{LAARA} \cite{Tripathi2026LAARA} provides a formal theoretical foundation for per-layer allocation using Fisher information and proves that uniform rank is suboptimal. Its implementation uses gradient warmup to estimate layer importance, followed by allocation. LAARA is the closest theoretical peer: it proves that per-layer allocation is optimal, while we investigate whether the allocation can be derived from $W_0$ alone.

\textbf{IGU-LoRA} \cite{Jiang2026IGULoRA} identifies a gradient bias in AdaLoRA's importance scores and proposes integrated gradient corrections. This motivates $W_0$-based approaches: if gradient-based signals are biased, structural signals may be more reliable.

\subsection{Calibration-Based Methods}

\textbf{IFCLoRA} \cite{Zhang2026IFCLoRA} is the closest methodological peer to our work. It computes per-layer importance from Information Flow Centrality (IFC) scores derived from a calibration forward pass. IFCLoRA demonstrates that pre-fine-tuning signals can predict rank need, motivating our hypothesis. The key difference: IFCLoRA requires calibration data and a forward pass; erank requires only the pretrained weight matrices. We view IFCLoRA as validating the principle and erank as testing its zero-data extreme.

\subsection{Structural Methods: $W_0$ as Initialization}

\textbf{PiSSA} \cite{Meng2024PiSSA} initializes LoRA matrices from the principal singular vectors of $W_0$, achieving faster convergence than random initialization. Rank is still set uniformly; the SVD of $W_0$ informs initialization, not rank.

\textbf{LoRA-XS} \cite{Balazy2024LoRAXS} freezes the full $W_0$ SVD as a structured matrix and trains only a small $r \times r$ core adapter. Like PiSSA, it uses $W_0$ structure for adapter design, not for rank selection.

These works establish that $W_0$'s singular structure contains adaptation-relevant information but do not test whether it predicts optimal per-layer rank.

\subsection{Intrinsic Dimensionality and Effective Rank}

\textbf{Aghajanyan et al.} \cite{Aghajanyan2021Intrinsic} showed that the intrinsic dimensionality of fine-tuning loss landscapes varies by layer, with $d_{90}$ differing substantially between attention and FFN layers. This is the foundational motivation for our hypothesis.

\textbf{Roy and Vetterli} \cite{RoyVetterli2007EffectiveRank} define effective rank $\text{erank}(A) = \exp(H(\sigma/\|\sigma\|_1))$, bounded $[1, \text{rank}(A)]$, scale-invariant, and providing wider dynamic range than stable rank or spectral norm ratios.

\subsection{Our Position}

We are the first to test whether $\text{erank}(W_0)$ --- a single number computed from pretrained weights before any fine-tuning --- correlates significantly with per-layer PARA oracle ranks. This places us in a novel position: using structural weight geometry as a zero-cost, zero-data rank predictor.
